% =================================================================
\section{3.3 나이브베이즈: 독립을 가정하고 차원의 저주를 피한다}
\begin{frame}{독립 가정으로 차원의 저주를 피하다}
  GDA는 저차원 연속값엔 잘 맞지만, 스팸처럼 \kb{어휘 크기(수만 개)}
  차원의 이산 데이터면 $\Sigma$ (어휘$\times$어휘) 추정이
  \textbf{감당이 안 됨}.
  \vskip0.5em
  \kb{나이브베이즈}는 과감한 단순화로 회피 ---
  \textbf{클래스 $y$ 가 주어지면 각 특징(단어)이 서로 독립}:
  \[ P(x|y) = \prod_{j=1}^n P(x_j|y) \]
  \vskip0.5em
  \begin{itemize}
    \item ``\kb{나이브(순진)}''이라는 이름의 뜻: 이 가정이
      \textbf{현실적으로 거의 항상 틀림} --- `무료'와 `당첨'은
      실제로 스팸에서 함께 나오는 경향.
    \item 그런데도 실전(특히 \textbf{텍스트 분류})에서 \textbf{놀랄
      만큼 잘 작동} --- 각 $P(x_j|y)$ 는 \textbf{단어 빈도만 세면}
      바로 추정 $\Rightarrow$ 어휘가 아무리 커도 파라미터가
      \textbf{어휘 크기에 선형}으로만 증가(GDA처럼 제곱이 아님).
  \end{itemize}
\end{frame}
